\ifdefined\gkver\else\def\gkver{student}\fi
\documentclass[\gkver]{gaokaozhenti}
\begin{document}

\chapter{1980年全国}

\begin{enumerate}
\item
%% number: 1
%% paperName: 1980年全国高考第 1 题 3 分
%% typeId: 1
%% type: 单选题
%% chapter: 力物体的平衡
%% point: 重力
%% method: 无
%% score: 3
%% degree: 850
%% duplicateId: 0
%% body:
月球表面上的重力加速度为地球表面上的重力加速度的 $\frac{1}{6}$，一个质量为 $600\Ukg$ 的飞行器在月球表面上\xzanswer{C}
\fourchoices[answer=C]
{质量是 $100\Ukg$，重力是 $5880\UN$}
{质量是 $100\Ukg$，重量是 $980\UN$}
{质量是 $600\Ukg$，重量是 $980\UN$}
{质量是 $600\Ukg$，重量是 $5880\UN$}

%% answer: C
%% memo:
\memoanswer{
因为质量是物体的固有属性，不随位置的改变而改变，所以质量仍然为 $600\Ukg$，故 AB 错误；由于在月球表面的重力是地球表面重力的 $\frac{1}{6}$，所以飞行器在月球表面的重力 $G_{\text{月}}=m\times\frac{1}{6}g_{\text{地}}=980\UN$，故 C 正确 D 错误。
}

\item
%% number: 2
%% paperName: 1980年全国高考第 2 题 3 分
%% typeId: 1
%% type: 单选题
%% chapter: 力物体的平衡
%% point: 受力分析
%% method: 无
%% score: 3
%% degree: 650
%% duplicateId: 0
%% body:
一架梯子斜靠在光滑的竖直墙上，下端放在水平的粗糙地面上。下面是梯子受力情况的简单描述，哪一句是正确的？梯子受到\xzanswer{C}
\fourchoices[answer=C]
{两个竖直的力，一个水平的力}
{一个竖直的力，两个水平的力}
{两个竖直的力，两个水平的力}
{三个竖直的力，两个水平的力}

%% answer: C
%% memo:
\memoanswer{
对梯子受力分析知，如图所示，梯子受到竖直向下的重力，竖直向上的地面支持力、水平方向墙面的弹力和地面的静摩擦力的作用，故 C 正确。

\onepicture[w=4.5cm]{figs/02_an.png}
}

\item
%% number: 3
%% paperName: 1980年全国高考第 3 题 3 分
%% typeId: 1
%% type: 单选题
%% chapter: 电路
%% point: 电容
%% method: 无
%% score: 3
%% degree: 850
%% duplicateId: 0
%% body:
一个平行板电容器，它的电容\xzanswer{B}
\fourchoices[answer=B]
{跟正对的面积成正比，跟两板间的距离成正比}
{跟正对的面积成正比，跟两板间的距离成反比}
{跟正对的面积成反比，跟两板间的距离成正比}
{跟正对的面积成反比，跟两板间的距离成反比}

%% answer: B
%% memo:
\memoanswer{
由平行板电容器的决定式 $C=\frac{\varepsilon S}{4\pi kd}$ 知，电容的大小跟正对的面积成正比，跟板间距成反比，故 B 正确。
}

\item
%% number: 4
%% paperName: 1980年全国高考第 4 题 3 分
%% typeId: 1
%% type: 单选题
%% chapter: 交流电
%% point: 交流电
%% method: 无
%% score: 3
%% degree: 850
%% duplicateId: 0
%% body:
把 $220\UV$ 的交流电压加在 $440\UO$ 的电阻上，在电阻上\xzanswer{A}
\fourchoices[answer=A]
{电压的有效值为 $220\UV$，电流的有效值为 $0.5\UA$}
{电压的最大值为 $220\UV$，电流的有效值为 $0.5\UA$}
{电压的有效值为 $220\UV$，电流的最大值为 $0.5\UA$}
{电压的最大值为 $220\UV$，电流的最大值为 $0.5\UA$}

%% answer: A
%% memo:
\memoanswer{
$220\UV$ 的交流电压指交流电压的有效值为 $220\UV$，电压的最大值为 $220\sqrt{2}\UV\approx311\UV$，故 BD 错误；将该交流电压加在 $440\UO$ 的电阻上，通过电阻的电流为 $I=\frac{220}{440}\UA=0.5\UA$，该电流的最大值为 $0.5\sqrt{2}\UA$，故 A 正确 C 错误。
}

\item
%% number: 5
%% paperName: 1980年全国高考第 5 题 3 分
%% typeId: 1
%% type: 单选题
%% chapter: 光学
%% point: 全反射
%% method: 无
%% score: 3
%% degree: 850
%% duplicateId: 0
%% body:
某媒质对空气的折射率为 1.414，一束光从媒质射向空气，入射角为 $60^{\circ}$。它的光路图是\xzanswer{D}
\fourchoices[answer=D, ispicture=true, h=2.2cm]
{\includesvg{figs/05a}}
{\includesvg{figs/05b}}
{\includesvg{figs/05c}}
{\includesvg{figs/05d}}

%% answer: D
%% memo:
\memoanswer{
光线从折射率为 $\sqrt{2}$ 的介质中射向空气，发生全反射的临界角为 $\sin C=\frac{1}{n}=\frac{\sqrt{2}}{2}=\sin 45^{\circ}$，由于入射角为 $60^{\circ}$，大于临界角，故会发生全反射，只有反射光线，故 D 正确。
}

\item
%% number: 6
%% paperName: 1980年全国高考第 6 题 3 分
%% typeId: 1
%% type: 单选题
%% chapter: 直线运动
%% point: 运动图象
%% method: 无
%% score: 3
%% degree: 650
%% duplicateId: 0
%% body:
物体竖直上抛后又落向地面，设向上的速度为正，它在整个运动过程中速度 $v$ 跟时间 $t$ 的关系是\xzanswer{B}
\fourchoices[answer=B, ispicture=true, h=2.8cm]
{\includesvg{figs/06a}}
{\includesvg{figs/06b}}
{\includesvg{figs/06c}}
{\includesvg{figs/06d}}

%% answer: B
%% memo:
\memoanswer{
A 图表示物体先向上匀减速运动，后向上匀加速运动，不符合物体的实际运动情况，故 A 错误；B 图表示物体先向上做匀减速直线运动，后向下做匀加速直线运动，加速度不变，故 B 正确；C 图表示物体先向下做匀减速运动，后向上做匀加速运动，不符合物体的实际运动情况，故 C 错误；D 图表示速度均为正值，表示物体速度方向不变，不符合物体的实际运动情况，故 D 错误。
}

\item
%% number: 7
%% paperName: 1980年全国高考第 7 题 3 分
%% typeId: 1
%% type: 单选题
%% chapter: 电场
%% point: 电势能电势差电场力做功
%% method: 无
%% score: 3
%% degree: 650
%% duplicateId: 0
%% body:
在电场中，A 点的电势高于 B 点的电势，\xzanswer{A}
\fourchoices[answer=A]
{把负电荷从 A 点移到 B 点，电场力作负功}
{把负电荷从 A 点移到 B 点，电场力作正功}
{把正电荷从 A 点移到 B 点，电场力作负功}
{把正电荷从 B 点移到 A 点，电场力作正功}

%% answer: A
%% memo:
\memoanswer{
因 A 点的电势高于 B 点的电势，由 $E_{p}=q\varphi$ 知，负电荷在 A 处电势能小于在 B 处电势能，把负电荷从 A 点移到 B 点，电势能增大，电场力做负功，故 A 正确 B 错误；与 A 的分析同理，由 $E_{p}=q\varphi$ 知，正电荷在 A 点的电势能大于在 B 的电势能，把正电荷从 A 点移到 B 点，电场力做正功，把正电荷从 B 点移到 A 点，电场力做负功，故 C 错误 D 错误。
}

\item
%% number: 8
%% paperName: 1980年全国高考第 8 题 3 分
%% typeId: 1
%% type: 单选题
%% chapter: 气体性质
%% point: 能量的转化及其方向性
%% method: 无
%% score: 3
%% degree: 650
%% duplicateId: 0
%% body:
把质量为 $m_{1}$ 的 $0\UdC$ 的冰放进质量为 $m_{2}$ 的 $100\UdC$ 的水中，最后得到温度是 $10\UdC$ 的水。如果容器吸热和热量损失均可忽略，那么 $m_{1}$ 和 $m_{2}$ 的关系是\xzanswer{D}
\fourchoices[answer=D]
{$m_{1}=9m_{2}$}
{$m_{1}=17m_{2}$}
{$6.1m_{1}=m_{2}$}
{$m_{1}=m_{2}$}

%% answer: D
%% memo:
\memoanswer{
水的相变潜热为 $\lambda=334\UkJkg$，质量为 $m_{1}$ 的冰吸热融化后再升温到 $10\UdC$ 吸收的热量等于质量为 $m_{2}$ 的水降温至 $10\UdC$ 放出的热量，即 $\lambda m_{1}+cm_{1}\Delta t_{1}=cm_{2}\Delta t_{2}$，其中 $\Delta t_{1}=10\UdC$，$\Delta t_{2}=90\UdC$，代入得 $334m_{1}+4.2m_{1}\times10=4.2m_{2}\times90$，化简得 $376m_{1}=378m_{2}$，所以 $m_{1}\approx m_{2}$，故 D 正确。
}

\item
%% number: 9
%% paperName: 1980年全国高考第 9 题 3 分
%% typeId: 1
%% type: 单选题
%% chapter: 气体性质
%% point: 热力学第一定律
%% method: 无
%% score: 3
%% degree: 650
%% duplicateId: 0
%% body:
一定质量的理想气体吸热膨胀，并保持压强不变，则它的内能增加，\xzanswer{C}
\fourchoices[answer=C]
{它吸收的热量等于内能的增量}
{它吸收的热量小于内能的增量}
{它吸收的热量大于内能的增量}
{它吸收的热量可以大于内能的增量，也可以小于内能的增量}

%% answer: C
%% memo:
\memoanswer{
一定质量的理想气体吸热膨胀，保持压强不变，它的内能增加，又气体对外做功，故吸收的热量大于内能的增加量，故 ABD 错误 C 正确。
}

\item
%% number: 10
%% paperName: 1980年全国高考第 10 题 3 分
%% typeId: 1
%% type: 单选题
%% chapter: 电磁感应
%% point: 楞次定律
%% method: 无
%% score: 3
%% degree: 450
%% duplicateId: 0
%% body:
如右图所示，一水平放置的矩形闭合线圈 abcd，在细长磁铁的 N 极附近竖直下落，保持 bc 边在纸外，ad 边在纸内，由图中的位置 Ⅰ 经过位置 Ⅱ 到位置 Ⅲ，位置 Ⅰ 和 Ⅲ 都很靠近 Ⅱ。在这个过程中，线圈中感生电流\xzanswer{A}
\onepicture[w=7cm]{\includesvg{figs/10}}
\fourchoices[answer=A]
{沿 abcd 流动}
{沿 dcba 流动}
{由 Ⅰ 到 Ⅱ 是沿 abcd 流动，由 Ⅱ 到 Ⅲ 是沿 dcba 流动}
{由 Ⅰ 到 Ⅱ 是沿 dcba 流动，由 Ⅱ 到 Ⅲ 是沿 abcd 流动}

%% answer: A
%% memo:
\memoanswer{
细长磁铁附近是非匀强磁场，由条形磁铁的磁场分布可知，线圈在位置 Ⅱ 时穿过矩形闭合线圈的磁通量为零。线圈从位置 Ⅰ 到位置 Ⅱ，从下向上穿过 abcd 的磁通量在减少，线圈从位置 Ⅱ 到位置 Ⅲ，从上向下穿过 abcd 的磁通量在增加，由楞次定律和右手螺旋定则可判知感应电流的方向是 abcd，故 A 正确。
}

\item
%% number: 11
%% paperName: 1980年全国高考第 11 题 4 分
%% typeId: 3
%% type: 填空题
%% chapter: 电路
%% point: 电动势
%% method: 无
%% score: 4
%% degree: 850
%% duplicateId: 0
%% body:
12 个相同的电池联接如图，每一个电池的电动势为 $1.5\UV$，内电阻为 $0.3\UO$，电池组的总电动势为 \tkanswer[1.2cm]{$6$} $\UV$，总内电阻为 \tkanswer[1.2cm]{$0.4$} $\UO$。
\onepicture[align=right, w=6.5cm]{figs/11.png}

%% answer: 6，0.4
\jdanswer{
$6\UV$，$0.4\UO$
}
%% memo:
\memoanswer{
每条支路 4 个电池串联，总电动势为 $6\UV$，三条支路并联电池组总电动势为单条支路电动势，故为 $6\UV$；同时由于单个电池内阻为 $0.3\UO$，故单条支路总内阻为 $1.2\UO$，三条支路并联，并联内阻为其 $\frac{1}{3}$，故为 $0.4\UO$。
}

\item
%% number: 12
%% paperName: 1980年全国高考第 12 题 4 分
%% typeId: 3
%% type: 填空题
%% chapter: 原子和原子核
%% point: 天然放射性
%% method: 无
%% score: 4
%% degree: 650
%% duplicateId: 0
%% body:
$\ce{^{238}_{92}U}$ 衰变成 $\ce{^{222}_{86}Rn}$，共发生了 \tkanswer{4} 次 $\alpha$ 衰变，\tkanswer{2} 次 $\beta$ 衰变。

%% answer: 4，2
\jdanswer{
4，2
}
%% memo:
\memoanswer{
核反应方程满足核电荷数和质量数守恒，$\ce{^{238}_{92}U}$ 衰变为 $\ce{^{222}_{86}Rn}$，设经过了 $m$ 次 $\alpha$ 衰变，$n$ 次 $\beta$ 衰变，由核电荷数和质量数守恒得 $92=86+2m-n$，$238=222+4m$，解得 $m=4$，$n=2$，故发生了 4 次 $\alpha$ 衰变，2 次 $\beta$ 衰变。
}

\item
%% number: 13
%% paperName: 1980年全国高考第 13 题 4 分
%% typeId: 3
%% type: 填空题
%% chapter: 光学
%% point: 光的电磁说
%% method: 无
%% score: 4
%% degree: 650
%% duplicateId: 0
%% body:
一束光自真空射入玻璃。已知光在真空中的波长为 $0.60\Uum$，传播速度为 $3.0\times10^{8}\Ums$，玻璃的折射率为 1.5。则光进入玻璃后的波长为 \tkanswer[2cm]{$4\times10^{-5}\Ucm$}，频率为 \tkanswer[2.4cm]{$5\times10^{14}$}$\UHz$。

%% answer: 4×10^-5，5×10^14
\jdanswer{
$4\times10^{-5}\Ucm$，$5\times10^{14}\UHz$
}
%% memo:
\memoanswer{
由 $n=\frac{c}{v}=\frac{\lambda_{0}}{\lambda}$ 得 $\lambda=\frac{\lambda_{0}}{n}=\frac{0.60\times10^{-6}}{1.5}=0.40\times10^{-6}=0.40\Uum$。
又 $c=\lambda_{0}\nu$，故 $\nu=\frac{c}{\lambda_{0}}=\frac{3\times10^{8}}{0.60\times10^{-6}}\UHz=5\times10^{14}\UHz$。
}

\item
%% number: 14
%% paperName: 1980年全国高考第 14 题 4 分
%% typeId: 3
%% type: 填空题
%% chapter: 功和能
%% point: 机械能守恒定律
%% method: 无
%% score: 4
%% degree: 650
%% duplicateId: 0
%% body:
一个小物体从光滑半球的顶点滑下，初速度很小，可以忽略不计，球半径为 $0.40\Um$，物体落地时速度的大小是 \tkanswer[1.5cm]{$2.8$} $\Ums$。
\onepicture[align=right, w=6cm]{\includesvg{figs/14}}

%% answer: 2.8
\jdanswer{
$2.8\Ums$
}
%% memo:
\memoanswer{
小球在运动过程中只有重力做功，由机械能守恒得 $mgR=\frac{1}{2}mv^{2}$，解得 $v=2\sqrt{2}\Ums\approx2.8\Ums$。
}

\item
%% number: 15
%% paperName: 1980年全国高考第 15 题 4 分
%% typeId: 3
%% type: 填空题
%% chapter: 气体性质
%% point: 玻意耳定律
%% method: 无
%% score: 4
%% degree: 850
%% duplicateId: 0
%% body:
在标准状态下，一细长而均匀的玻璃管，上端开口，一段长度为 $38\Ucm$ 的水银柱，把一定量的空气封闭在管内。当玻璃管跟竖直方向成 $60^{\circ}$ 时，管内空气柱的长度为 $60\Ucm$。如果使管竖立，在管内空气达到平衡状态后，这段空气柱的长度是 \tkanswer[1.5cm]{$50$} $\Ucm$。
\onepicture[align=right, w=5cm]{figs/15.png}

%% answer: 50
\jdanswer{
$50\Ucm$
}
%% memo:
\memoanswer{
设倾斜时管内空气的压强为 $p_{1}$，竖直时管内空气的压强为 $p_{2}$，玻璃管的横截面积为 $S$，水银的密度为 $\rho$，由液柱受力平衡得

$p_{1}S=p_{0}S+\rho gh\sin30^{\circ}S$，

$p_{2}S=p_{0}S+\rho ghS$，

其中，$p_{0}$ 为标准大气压，$p_{0}=76\UcmHg$，水银柱的长度为 $h=38\Ucm$。

对管内空气，由玻意耳定律得

$p_{1}lS=p_{2}l'S$，

其中 $l$、$l'$ 分别表示倾斜和竖直状态空气柱的长度，$l=60\Ucm$。

联立各式并代入数据解得 $l'=50\Ucm$。
}

\item
%% number: 16
%% paperName: 1980年全国高考第 16 题 10 分
%% typeId: 8
%% type: 简答题
%% chapter: 光学
%% point: 透镜
%% method: 无
%% score: 10
%% degree: 450
%% duplicateId: 0
%% body:
在薄凸透镜成像中，设 $u$ 表示物距，$v$ 表示像距，$f$ 表示透镜的焦距。试证明薄凸透镜成像的公式为 $\frac{1}{u}+\frac{1}{v}=\frac{1}{f}$（可只证明成实像的情况）。再根据上式找出像距 $v$ 为正负值的条件，并指出它跟像的虚实的关系。

%% answer: （1）绘出凸透镜成像的光路图；（2）见解析；（3）$u>f$ 时 $v$ 为正值成实像，$u<f$ 时 $v$ 为负值成虚像
\jdanswer{
\begin{enumerate}
	\item
	绘出凸透镜成像的光路图。
	\item
	证明见解析。
	\item
	当 $u>f$ 时，$v$ 为正值，所成的是实像；当 $u<f$ 时，$v$ 为负值，所成的是虚像。
\end{enumerate}
}
%% memo:
\memoanswer{
（1）绘出凸透镜成像的光路图。

\onepicture[w=8cm]{figs/16_an.png}

（2）从光路图中可以看出

$\Delta ABO\sim\Delta A'B'O$，得 $\frac{AB}{A'B'}=\frac{BO}{B'O}$；①

$\Delta O'OF'\sim\Delta A'B'F'$，得 $\frac{O'O}{A'B'}=\frac{O'F'}{B'F'}$；②

由于 $AB=O'O$，因此

$\frac{BO}{B'O}=\frac{OF'}{B'F'}$；

由于 $BO=u$，$B'O=v$，$OF'=f$，$B'F'=v-f$，因此

$\frac{u}{v}=\frac{f}{v-f}$；③

化简后可得

$\frac{1}{u}+\frac{1}{v}=\frac{1}{f}$。④

这就是所要证明的薄凸透镜成像的公式。

（3）从上式可得：

$\frac{1}{v}=\frac{1}{f}-\frac{1}{u}=\frac{u-f}{fu}$。

由于 $fu$ 总是正值，如果 $u>f$，即物体位于焦点之外，则 $v$ 为正值，所成的是实像；如果 $u<f$，即物体位于焦点之内，则 $v$ 为负值，所成的是虚像。

评分说明：全题 10 分。（1）2 分，（2）4 分，（3）4 分。

（2）中，只找出相似三角形并列出①、②式的给 1 分，进一步导出③式的，再给 2 分。

（3）中，只找出 $v$ 为正负值的条件的，给 2 分；没有找出上述条件，只是写出 $v$ 为正值时成实像、$v$ 为负值时成虚像的，给 2 分。
}

\item
%% number: 17
%% paperName: 1980年全国高考第 17 题 12 分
%% typeId: 6
%% type: 计算题
%% chapter: 磁场
%% point: 带电粒子在复合场中的运动
%% method: 无
%% score: 12
%% degree: 550
%% duplicateId: 0
%% body:
如右图所示，一束具有各种速率的带一个基本正电荷的两种铜离子，质量数分别为 63 和 65，水平地经小孔 S 进入有匀强电场和匀强磁场的区域。电场 $E$ 的方向向下，磁场 $B$ 的方向垂直纸面向里。只有那些路径不发生偏折的离子才能通过另一个小孔 $S'$。为了把从 $S'$ 射出的两种铜离子分开，再让它们进入另一方向垂直纸面向外的匀强磁场 $B'$ 中，使两种离子分别沿不同半径的圆形轨道运动。试分别求出两种离子的轨道半径。（应明确说明演算过程的物理上的根据）

已知：$E=1.00\times10^{5}\UVm$，$B=0.4\UT$，$B'=0.50\UT$，基本电荷 $e=1.60\times10^{-19}\UC$，质量数为 63 的铜原子的质量 $m_{1}=63\times1.66\times10^{-27}\Ukg$，质量数为 65 的铜原子的质量 $m_{2}=65\times1.66\times10^{-27}\Ukg$。
\onepicture[align=right, w=5.5cm]{figs/17.png}

%% answer: $R_{1}=0.33$ m，$R_{2}=0.34$ m
\jdanswer{
$R_{1}\approx0.33\Um$，$R_{2}\approx0.34\Um$
}
%% memo:
\memoanswer{
（1）设铜离子的电量为 $e$，以速度 $v$ 进入小孔 S 后，受到的力有电场力 $F_{1}=eE$，方向向下，洛伦兹力 $F_{2}=evB$，方向向上，重力可忽略不计，只有当 $F_{1}=F_{2}$ 时，铜离子才能匀速无偏折地穿出小孔 $S'$。因此，从小孔 $S'$ 穿出的铜离子必须满足的条件是：

$eE=evB$。①

这就是说，只有速度 $v=\frac{E}{B}$ 的铜离子能穿出小孔 $S'$。

（2）铜离子进入磁场 $B'$ 后，受到洛伦兹力 $F=evB'$，重力仍可忽略不计。$F$ 跟 $v$ 垂直并为一恒量，因此铜离子在磁场 $B'$ 内将作匀速圆周运动，$F$ 就是这种圆周运动的向心力，设铜 63 离子和铜 65 离子运动轨迹的半径分别为 $R_{1}$ 和 $R_{2}$，那么，

$evB'=m_{1}\frac{v^{2}}{R_{1}}$；②

$evB'=m_{2}\frac{v^{2}}{R_{2}}$。③

（3）由①、②两式可得：

$R_{1}=\frac{m_{1}E}{eBB'}$；④

由①、③两式可得：

$R_{2}=\frac{m_{2}E}{eBB'}$。⑤

代入数值进行计算，

$R_{1}=\frac{63\times1.66\times10^{-27}\times1.00\times10^{5}}{1.60\times10^{-19}\times0.40\times0.50}\Um=0.33\Um$，

$R_{2}=\frac{65\times1.66\times10^{-27}\times1.00\times10^{5}}{1.60\times10^{-19}\times0.40\times0.50}\Um=0.34\Um$。

评分说明：全题 12 分。（1）4 分，（2）4 分，（3）4 分。

（1）中，能正确写出 $F_{1}$ 和 $F_{2}$ 的表达式并指明方向的，各给 1 分。未说明重力可忽略不计的不扣分；能进一步指出只有当 $F_{1}$ 跟 $F_{2}$ 平衡时，铜离子才能穿出 $S'$，并列出①式的，再给 2 分；没有任何说明直接列出①式的，只给 2 分。

（2）中，只写出 $F$ 的表达式并指出 $F$ 跟 $v$ 垂直的，给 1 分；进一步说明 $F$ 就是向心力的，再给 1 分；进一步列出②、③两式的，再各给 1 分，没有任何说明直接列出②、③两式的，给 3 分。

（3）中，只导出④、⑤两式的，各给 1 分；单纯数字计算错误的，扣 1 分；答案中未写单位或单位错误的，扣 1 分。
}

\item
%% number: 18
%% paperName: 1980年全国高考第 18 题 14 分
%% typeId: 4
%% type: 实验题
%% chapter: 电路
%% point: 电路实验
%% method: 无
%% score: 14
%% degree: 450
%% duplicateId: 0
%% body:
右图表示用惠斯登电桥来测量未知电阻 $R_{x}$。图中 $R$ 是已知电阻；K 是电键；G 是电流计；AC 是一条粗细均匀的长直电阻丝；D 是滑动触头，按下时就使电流计的一端与电阻丝接通；L 是米尺。
\begin{enumerate}
	\item
	简要说明测量 $R_{x}$ 的实验步骤（其中要写出计算 $R_{x}$ 的公式）。
	\item
	如果滑动触头 D 在从 A 向 C 移动的整个过程中，每次按下 D 时，流过 G 的电流总是比前一次增大，已知 AC 间的电阻丝是导通的，那么，电路可能在什么地方断路了？说明理由。（分析时可认为电池内电阻和电阻 $R'$ 均可忽略）。
\end{enumerate}
\onepicture[align=right, w=6cm]{figs/18.png}

%% answer: （1）见解析；（2）BC 间发生了断路
\jdanswer{
\begin{enumerate}
	\item
	①按下电键 K，把滑动触头放在 AC 中点附近，按下 D，观察电流计 G 中指针偏转方向；②向左或向右移动 D，直到按下 D 时 G 的指针不偏转为止；③用米尺量出此时 AD、DC 的长度 $l_{1}$ 和 $l_{2}$；④根据公式 $R_{x}=\frac{l_{2}}{l_{1}}R$ 计算出 $R_{x}$ 的值。
	\item
	BC 间发生了断路。
\end{enumerate}
}
%% memo:
\memoanswer{
（1）实验步骤：①按下电键 K，把滑动触头放在 AC 中点附近，按下 D，观察电流计 G 中指针偏转方向；②向左或向右移动 D，直到按下 D 时 G 的指针不偏转为止；③用米尺量出此时 AD、DC 的长度 $l_{1}$ 和 $l_{2}$；④根据公式 $R_{x}=\frac{l_{2}}{l_{1}}R$ 计算出 $R_{x}$ 的值。

（2）BC 间发生了断路。

理由：此时的等效电路如图所示，D 向右移动 $R_{1}$ 增大，AD 间的电阻 $R_{\mathrm{AD}}=\frac{R_{1}R_{3}}{R_{1}+R_{3}}$ 也增大，其中 $R_{3}=R+R_{G}$；$R_{G}$ 是电流计内电阻，同时 $R_{2}$ 减小。于是加在 AD 间的电压 $U_{\mathrm{AD}}=\frac{R_{\mathrm{AD}}E}{R_{\mathrm{AD}}+R_{2}}=\frac{E}{1+\frac{R_{2}}{R_{\mathrm{AD}}}}$

\onepicture[w=6cm]{figs/18_an.png}

也增大，这就使流过电流计 G 的电流增大。

评分说明：全题 14 分。（1）8 分，（2）6 分。

（1）中，步骤①2 分，未写“按下 K”，或未写“按下 D，观察 G 中指针偏转方向”的，各扣 1 分；步骤②3 分，在按下 D 时未写“直到 G 的指针不偏转”的不给分，漏写按下 D 的，扣 1 分；步骤③1 分；步骤④2 分，未写 $R_{x}=\frac{l_{2}}{l_{1}}R$ 的，不给分。$l_{1}$，$l_{2}$ 用 AD、DC 表示也可以。如果步骤①的全部或部分内容与步骤②合并写出，不因此而扣分。

（2）中，写出断路位置，给 2 分，说明理由，给 4 分。理由的说明中，只写出 $R_{\mathrm{AD}}$ 随 D 的右移而增大的，给 2 分。用分压原理说明了随 D 的右移 AD 间的电压增大，因而流过 G 的电流增大的，给 4 分。说明理由中，只要基本点都已说到，无概念错误，但叙述层次不清、文句不够通顺的，这次考试暂不扣分，明显的概念错误可酌情扣分。
}

\item
%% number: 19
%% paperName: 1980年全国高考第 19 题 14 分
%% typeId: 6
%% type: 计算题
%% chapter: 运动和力
%% point: 牛顿定律的应用
%% method: 分情况讨论
%% score: 14
%% degree: 450
%% duplicateId: 0
%% body:
有两个物体，质量分别为 $m_{1}$ 和 $m_{2}$。$m_{1}$ 原来静止，$m_{2}$ 以速度 $v$ 向右运动，如图所示。
\begin{enumerate}
	\item
	它们同时开始受到向右的大小相同的恒力 $F$，在 $m_{1}<m_{2}$，$m_{1}=m_{2}$，$m_{1}>m_{2}$ 三种情况下，它们能否达到相同的速度（矢量）？
	\item
	试列出它们速度的表达式，并根据此式分别进行讨论，讨论中要注意说明理由。
	\item
	如果它们受到的恒力 $F$ 的方向都跟 $v$ 垂直，它们能否达到相同的速度（矢量）？为什么？
\end{enumerate}
\onepicture[align=right, w=5cm]{\includesvg{figs/19}}

%% answer: （1）$v_{1}=\frac{F}{m_{1}}t$，$v_{2}=v+\frac{F}{m_{2}}t$；（2）见解析；（3）不能达到相同的速度
\jdanswer{
\begin{enumerate}
	\item
	$v_{1}=\frac{F}{m_{1}}t$，$v_{2}=v+\frac{F}{m_{2}}t$。
	\item
	见解析。
	\item
	不能达到相同的速度。
\end{enumerate}
}
%% memo:
\memoanswer{
（1）设受力后 $m_{1}$ 的加速度为 $a_{1}$，$m_{2}$ 的加速度为 $a_{2}$，受力后某一时刻 $t$，$m_{1}$ 的速度为 $v_{1}$，$m_{2}$ 的速度为 $v_{2}$，那么：

$a_{1}=\frac{F}{m_{1}}$，$a_{2}=\frac{F}{m_{2}}$；

$v_{1}=a_{1}t=\frac{F}{m_{1}}t$；①

$v_{2}=v+a_{2}t=v+\frac{F}{m_{2}}t$。②

（2）受力后，$m_{1}$ 做初速为零的匀加速运动，$m_{2}$ 做有一定初速度的匀加速运动，它们的加速度和速度的方向都是向右的。

$m_{1}<m_{2}$ 时，由于 $a_{1}>a_{2}$，$m_{1}$ 的速度增加得比 $m_{2}$ 的快，虽然 $m_{2}$ 已有一定初速度，它们仍可在某一时刻达到相同的速度。

$m_{1}=m_{2}$ 时，由于 $a_{1}=a_{2}$，它们的速度增加得一样快，$m_{2}$ 已有一初速度 $v$，因此 $m_{1}$ 的速度将总是比 $m_{2}$ 的速度小 $v$，它们不可能达到相同的速度。

$m_{1}>m_{2}$ 时，由于 $a_{1}<a_{2}$，$m_{1}$ 的速度增加得比 $m_{2}$ 的慢，$m_{2}$ 已有一初速度，因此 $m_{1}$ 的速度将越来越小于 $m_{2}$ 的速度，它们也不可能达到相同的速度。

（3）如果 $F$ 跟 $v$ 垂直，那么，$F$ 的作用只是使 $m_{1}$ 和 $m_{2}$ 在垂直于 $v$ 的方向上的速度增加，而对它们在 $v$ 方向的即向右的速度没有影响。因此 $m_{1}$ 将始终没有向右的速度分量，而 $m_{2}$ 将在向右的方向上始终保持速度 $v$。这样，在任何时刻 $m_{1}$ 和 $m_{2}$ 的速度方向都不会相同，因此它们不可能达到相同的速度。

评分说明：全题 14 分。（1）2 分，（2）7 分，（3）5 分。

（1）未写出 $v_{1}=\frac{F}{m_{1}}t$ 或 $v_{2}=v+a_{2}t=v+\frac{F}{m_{2}}t$ 的，各扣 1 分。

（2）中，正确说明了 $m_{1}$ 和 $m_{2}$ 受力后的运动情况的，给 3 分，但说明中未明确指出加速度和速度方向的，扣 1 分；对三种情况的讨论，每个正确解答给 2 分；对 $m_{1}=m_{2}$ 的讨论中未说到 $v_{1}$ 总比 $v_{2}$ 小 $v$ 的，对 $m_{1}>m_{2}$ 的讨论中未说到 $v_{1}$ 比 $v_{2}$ 越来越小的，均不扣分。

（3）中，未指出 $F$ 对 $v$ 方向上的速度没有影响的，扣 2 分；未指出 $v_{1}$ 和 $v_{2}$ 的方向不可能相同的，扣 3 分。

考生也可以根据①、②两式，设 $v_{1}=v_{2}$，求出它们速度相等的时间 $t=\frac{mv_{1}v_{2}}{m_{2}-m_{1}}$ 后进行讨论，从 $t$ 有正解、无解或负解来说明 $m_{1}$ 和 $m_{2}$ 能否达到相同的速度，只要说理正确，应同样给分。答卷中只要基本点都已说到，无概念错误，但叙述层次不清、文句不够通顺的，这次考试暂不扣分，明显的概念错误可酌情扣分。
}

\end{enumerate}

\end{document}
